\documentclass[11pt]{article}

\usepackage[a4paper,margin=1in]{geometry}
\usepackage[T1]{fontenc}
\usepackage[utf8]{inputenc}
\usepackage{lmodern}
\usepackage{amsmath,amssymb,mathtools}
\usepackage{booktabs}
\usepackage{xcolor}
\usepackage[colorlinks=true,linkcolor=blue!60!black,urlcolor=blue!60!black]{hyperref}

\newcommand{\GeV}{\mathrm{GeV}}
\newcommand{\MeV}{\mathrm{MeV}}
\newcommand{\keV}{\mathrm{keV}}
\newcommand{\ii}{\mathrm{i}}

\title{$D_{s1}$ normalized-current $AA/AB/BB$ two-point QCD sum rule}
\author{Calculation note}
\date{\today}

\begin{document}
\maketitle

\begin{abstract}
This note replaces the former external-$f_1$/overlap normalization of the
$D_{s1}(2460)$ and $D_{s1}(2536)$ currents.  One normalized-current
$AA/AB/BB$ two-point QCD sum-rule matrix determines the mixing angle and both
physical residues.  The calculation is implemented independently in Python
and Mathematica.  At the stated LO+$d=3$+$d=5$ truncation the central
Mathematica and Python results agree below $10^{-7}$.
\end{abstract}

\section{Scope and the local-condensate issue}

This is an ordinary two-point SVZ sum rule.  Therefore the local vacuum
condensates $\langle\bar ss\rangle$ and
$\langle\bar s g_s\sigma\!\cdot\!G s\rangle$ are part of its OPE.  They are
not inserted into the external-photon transition LCSR, which follows the
Rohrwild organization in terms of non-local photon matrix elements and photon
distribution amplitudes.  The two statements concern different correlators
and are fully compatible.

The present common matrix truncation contains
\begin{itemize}
  \item exact-$m_s$ leading-order perturbative spectral densities;
  \item the local strange condensate through $O(m_s^2)$;
  \item the local dimension-five mixed condensate.
\end{itemize}
The gluon condensate and perturbative $\alpha_s$ corrections are not included.
Thus ``complete'' below means a complete $2\times2$ current matrix at this
explicit common truncation, not an all-dimensional or all-orders OPE.

\section{Currents, correlator, and rotation}

The normalized basis currents are
\begin{align}
J_A^\mu&=\bar s\gamma^\mu\gamma_5c,\qquad
J_B^\mu=\frac{\ii}{m_c+m_s}\bar s\sigma^{\mu\nu}p'_\nu\gamma_5c,
\end{align}
where $p'$ is the initial-state momentum.  The physical currents are
\begin{align}
J_1^\mu&=\sin\theta\,J_A^\mu+\cos\theta\,J_B^\mu,\\
J_2^\mu&=\cos\theta\,J_A^\mu-\sin\theta\,J_B^\mu.
\end{align}
For $i,j\in\{A,B\}$,
\begin{equation}
\Pi_{ij}^{\mu\nu}(p')=\ii\int d^4x\,e^{\ii p'\cdot x}
\langle0|T\{J_i^\mu(x)J_j^{\nu\dagger}(0)\}|0\rangle
=\left(-g^{\mu\nu}+\frac{p'^\mu p'^\nu}{p'^2}\right)\Pi_{ij}(p'^2)+\cdots.
\end{equation}
Writing
\begin{equation}
R(\theta)=\begin{pmatrix}\sin\theta&\cos\theta\\
\cos\theta&-\sin\theta\end{pmatrix},\qquad
\Pi^{\rm phys}=R\Pi^{\rm basis}R^T,
\end{equation}
gives
\begin{align}
\Pi_{11}&=\sin^2\theta\,\Pi_{AA}+\cos^2\theta\,\Pi_{BB}
+2\sin\theta\cos\theta\,\Pi_{AB},\label{eq:pi11}\\
\Pi_{22}&=\cos^2\theta\,\Pi_{AA}+\sin^2\theta\,\Pi_{BB}
-2\sin\theta\cos\theta\,\Pi_{AB},\\
\Pi_{12}&=\sin\theta\cos\theta(\Pi_{AA}-\Pi_{BB})
+(\cos^2\theta-\sin^2\theta)\Pi_{AB}.
\end{align}
Equation~\eqref{eq:pi11} follows directly from the stated definition of
$J_1$.  The early draft had the $AA$ and $BB$ coefficients interchanged in
this equation; that version is incorrect.

\section{Perturbative spectral matrix}

Let $d=m_c+m_s$ and
\begin{equation}
\lambda(s)=[s-(m_c+m_s)^2][s-(m_c-m_s)^2].
\end{equation}
The exact-mass LO densities are
\begin{align}
\rho_{AA}(s)&=\frac{\sqrt\lambda}{8\pi^2}\left[
2-\frac{m_c^2+m_s^2+6m_cm_s}{s}
-\frac{(m_c^2-m_s^2)^2}{s^2}\right],\\
\rho_{AB}(s)&=\rho_{BA}(s)=\frac{\sqrt\lambda}{8\pi^2}
\frac{3(m_s-m_c)(d^2-s)}{ds},\\
\rho_{BB}(s)&=\frac{\sqrt\lambda}{8\pi^2}
\frac{-(d^2-s)(2m_s^2-4m_cm_s+2m_c^2+s)}{d^2s}.
\end{align}
The $AA$ density is algebraically identical to the Colangelo axial-current
two-point density.  The equality and $AB=BA$ are exact FeynCalc regression
checks, not numerical observations.

\section{Local Borel matrix}

With $C_3=\langle\bar ss\rangle$ and
$C_5=\langle\bar s g_s\sigma\!\cdot\!G s\rangle$, the continuum-subtracted
Borel matrix is
\begin{equation}
\widehat\Pi(M^2,s_0)=
\int_{d^2}^{s_0}ds\,e^{-s/M^2}\rho(s)
+e^{-m_c^2/M^2}\left[C_3(\mathcal Q_0+\mathcal Q_1+\mathcal Q_2)
+C_5\mathcal Q_5\right].
\end{equation}
The local matrices are
\begin{align}
\mathcal Q_0&=
\begin{pmatrix}m_c&m_c^2/d\\m_c^2/d&m_c^3/d^2\end{pmatrix},\\
\mathcal Q_1&=
\begin{pmatrix}
-\dfrac{m_sm_c^2}{2M^2}&
\dfrac{m_sm_c(1-m_c^2/M^2)}{2d}\\
\dfrac{m_sm_c(1-m_c^2/M^2)}{2d}&
\dfrac{m_sm_c^2(1-m_c^2/(2M^2))}{d^2}
\end{pmatrix},\\
\mathcal Q_2&=
\begin{pmatrix}
\dfrac{m_s^2m_c^3}{2M^4}&
\dfrac{m_s^2(m_c^4/M^4-m_c^2/M^2-1)}{2d}\\
\dfrac{m_s^2(m_c^4/M^4-m_c^2/M^2-1)}{2d}&
\dfrac{m_s^2(m_c^5/(2M^4)-m_c^3/M^2)}{d^2}
\end{pmatrix},\\
\mathcal Q_5&=
\begin{pmatrix}
-\dfrac{m_c^3}{4M^4}&
-\dfrac{m_c^4/M^4-m_c^2/M^2-1}{4d}\\
-\dfrac{m_c^4/M^4-m_c^2/M^2-1}{4d}&
-\dfrac{m_c^5/(4M^4)-m_c^3/(2M^2)}{d^2}
\end{pmatrix}.
\end{align}
The complete $AA$ local entry reduces exactly to
\begin{equation}
e^{-m_c^2/M^2}\left\{
C_3\left[m_c-\frac{m_c^2m_s}{2M^2}
+\frac{m_c^3m_s^2}{2M^4}\right]
-C_5\frac{m_c^3}{4M^4}\right\},
\end{equation}
which is the corresponding Colangelo expression.

\section{Fixed angle, thresholds, and residues}

The nominal angle is taken from our dedicated mixing-angle study,
\begin{equation}
\theta_{D_s}=26.6\pm0.6^\circ.
\end{equation}
The angle obtained from $\widehat\Pi_{12}=0$ in the present
LO+$d=3$+$d=5$ matrix is retained only as a truncation diagnostic.  We fit a
separate threshold $s_0^{(i)}$ to each measured pole and use
\begin{align}
m_{i,\mathrm{eff}}^2&=-\frac{d}{d(1/M^2)}
\ln\widehat\Pi_{ii}(M^2,s_0^{(i)};\theta),\\
f_i^2m_i^2e^{-m_i^2/M^2}
&=\widehat\Pi_{ii}(M^2,s_0^{(i)};\theta),\qquad i=1,2.
\end{align}

In 2000 input samples, 1714 pass all diagnostics.  The result is
\begin{align}
\theta_{D_s}&=26.595[25.996,27.199]^\circ,\\
\theta_{\rm matrix}&=30.632[29.942,31.286]^\circ,\\
f_1&=0.4046[0.3799,0.4291]~\GeV,&
f_2&=0.1677[0.1594,0.1756]~\GeV.
\end{align}
The normalized off-diagonal residual after projecting at the external angle is
$0.183[0.142,0.225]$.  It is an OPE-truncation diagnostic, not an additional
statistical angle error.  Neither residue is externally anchored.

At $M^2=2.35~\GeV^2$, the fixed-angle Python regression point is
\begin{equation}
\theta=26.6^\circ,\quad \theta_{\rm matrix}=30.58596159^\circ,\quad
s_0^{(1)}=7.925~\GeV^2,\quad s_0^{(2)}=9.975~\GeV^2,
\end{equation}
\begin{equation}
f_1=0.4005893052~\GeV,\qquad f_2=0.1666703674~\GeV.
\end{equation}
The independent Mathematica regression gives
\begin{equation}
f_1=0.4005892785~\GeV,\qquad
f_2=0.1666703614~\GeV .
\end{equation}
The differences from Python are $-2.67\times10^{-8}~\GeV$ and
$-6.02\times10^{-9}~\GeV$, respectively; the matrix-angle and normalized
$\Pi_{12}$ diagnostics agree to better than $10^{-7}$.

\section{Propagation, interference, and final charm widths}

The physical couplings are
\begin{align}
G_{2460}&=\frac{\sin\theta\,f_AG_A+\cos\theta\,f_BG_B}{f_1},\\
G_{2536}&=\frac{\cos\theta\,f_AG_A-\sin\theta\,f_BG_B}{f_2}.
\end{align}
For the preferred lattice photon normalization, the final-window Monte Carlo
gives
\begin{align}
G_{2460}&=-0.3569[-0.3929,-0.3237]~\GeV^{-1},&
\Gamma_{2460}&=26.71[22.01,32.44]~\keV,\\
G_{2536}&=-0.07305[-0.10434,-0.03836]~\GeV^{-1},&
\Gamma_{2536}&=1.656[0.4567,3.378]~\keV.
\end{align}
The low-state pieces are constructive, whereas the high-state pieces are
destructive throughout the physical first quadrant.  At the midpoint of the
selected transition window the benchmark widths are
\begin{center}
\begin{tabular}{lccc}
\toprule
Angle input & $\theta$ & $\Gamma_{2460}$ (keV) & $\Gamma_{2536}$ (keV)\\
\midrule
previous study & $26.6^\circ$ & $27.60$ & $1.623$\\
truncated-matrix diagnostic & $30.61^\circ$ & $27.71$ & $0.043$\\
HQET benchmark & $35.3^\circ$ & $27.64$ & $1.120$\\
\bottomrule
\end{tabular}
\end{center}
The high-state zero is near $31.37^\circ$ at this representative point.  The
interference mechanism does not change; the relative weights move the result
toward or away from the cancellation.

\section{Projection onto Wang's pure-AA channel}

Wang's arXiv v4 result is
$f_{D_{s1}}^{(J_A)}=0.345\pm0.017~\GeV$; the earlier
$0.245\pm0.017~\GeV$ entry was explicitly corrected.  Reproducing this
one-pole number is neither necessary nor sufficient to validate the mixed
residues.  The inverse current rotation is
\begin{equation}
J_A=\sin\theta\,J_1+\cos\theta\,J_2,
\end{equation}
so our physical poles project back to
\begin{equation}
\widehat\Pi_{AA}^{(2\,{\rm pole})}
=\sin^2\theta\,f_1^2M_1^2e^{-M_1^2/M^2}
+\cos^2\theta\,f_2^2M_2^2e^{-M_2^2/M^2}.
\end{equation}
The corresponding one-pole-equivalent diagnostic at the lower mass is
\begin{equation}
f_{A,\rm eq}^2(M^2)=
\frac{e^{M_1^2/M^2}}{M_1^2}\widehat\Pi_{AA}^{(2\,{\rm pole})}
\end{equation}
and gives $0.2303[0.2194,0.2418]~\GeV$.  Separately, in Wang's Borel and
threshold window our direct pure-$AA$ LO+$d=3$+$d=5$ OPE gives
\begin{equation}
f_A^{\rm ours}=0.2686[0.2607,0.2768]~\GeV,
\end{equation}
which does not reproduce Wang's absolute normalization.  These three numbers
use different pole models and OPE truncations.  They are useful normalization
benchmarks, not a proof or disproof of the mixed-current calculation.  Neither
$f_1$ nor $f_2$ should be equated to Wang's number.  The individual projected
overlaps are
\begin{align}
\sin\theta\,f_1&=0.1812[0.1692,0.1930]~\GeV,\\
\cos\theta\,f_2&=0.1498[0.1424,0.1573]~\GeV,
\end{align}
These are projection diagnostics only.

\section{Executable record}

\begin{itemize}
  \item Python matrix and Monte Carlo: \path{../scripts/twopoint_ds1_matrix_sumrule.py}.
  \item Spectral trace audit: \path{../scripts/twopoint_dirac_spectral_audit.wl}.
  \item Local trace audit: \path{../scripts/twopoint_local_condensate_audit.wl}.
  \item Numerical Mathematica regression:
    \path{../scripts/mathematica_twopoint_ds1_matrix_check.wl}.
  \item Cell-by-cell notebook:
    \path{../notebooks/Ds1_AA_AB_BB_two_point_sumrule.nb}.
\end{itemize}

The former $f_1=0.345~\GeV$, $f_2\simeq0.38~\GeV$ overlap model is preserved
only as a legacy comparison and is not the current paper normalization.

\end{document}
